\textbf{1:}\\ \noindent
\vspace{0.1cm}

Popular method for post-hoc ensemble selection: Greedy ensemble selection (GES\footnote{Caruana, R., Niculescu-Mizil, A., Crew, G., \& Ksikes, A. (2004, July). Ensemble selection from libraries of models. In Proceedings of the twenty-first International Conference on Machine Learning (p. 18). \url{https://dl.acm.org/doi/pdf/10.1145/1015330.1015432}})
\begin{enumerate}
    \item Start with the empty ensemble
    \item Add the base model that maximizes the ensembles performance on a hillclimb (validation) set
    \item Repeat step 2 until a given number of iterations are reached
    \item Return the ensemble from the set of generated ensembles that has maximum performance on the validation set
\end{enumerate}
Note that ensemble selection effectively optimizes the weights on an integer domain (i.e., base models are selected one or multiple times or not at all)

Write your own GES algorithm constructing an ensemble from a set of five base models obtained via fitting: DT, SVM, KNN, Elastic-Net, Logistic Regression.
Use the adult data set from OpenML.
Construct a 60/20/20 train/validation/test split manually (stratified for the target) and use the AUC ROC as performance metric.
Train each base learner on the train split to generate the base models and use the validation split to perform GES for 10 iterations. Which models are selected and how often are they selected?
What is the test performance of the final ensemble?
How does it compare to the test performance of each base model?
